\section{Desarrollo incremental y gestión del entorno}

\subsection{Importancia del desarrollo incremental}

Sobre la base de un feature list, se solicita al Coding Agent que procese solo una funcionalidad a la vez (one feature at a time). Este método incremental ha demostrado ser clave para resolver el problema de los agentes que tratan demasiadas cosas a la vez.

Los beneficios del desarrollo incremental incluyen:
\begin{itemize}
    \item Objetivos claros y alcance controlado en cada sesión
    \item Incluso si una sesión falla, las pérdidas se limitan al ámbito de una sola funcionalidad
    \item Se conserva el registro de implementación de cada funcionalidad mediante git commit, lo que permite realizar rollback
    \item El agente de la siguiente sesión puede localizar rápidamente el progreso actual
\end{itemize}

\subsection{Tareas de limpieza al finalizar la sesión}

Antes de terminar cada sesión del Coding Agent, se deben ejecutar las siguientes operaciones de limpieza:

\begin{enumerate}
    \item Comprobar los cambios de código en git, adjuntando un mensaje descriptivo del commit
    \item Escribir un resumen del trabajo realizado en esta sesión en \texttt{claude-progress.txt}
    \item Actualizar el estado de las funcionalidades completadas en \texttt{feature\_list.json}
    \item Asegurar que el código se encuentre en un estado compilable y ejecutable
\end{enumerate}

\begin{knowledgebox}{Git como red de seguridad}
Git no es solo una herramienta de control de versiones; en este sistema también desempeña el papel de "red de seguridad". Cuando las modificaciones del código del agente introducen errores o provocan regresiones, el agente puede utilizar \texttt{git revert} para volver al estado de trabajo anterior, en lugar de dedicar mucho tiempo a la depuración. Esta estrategia de "rollback primero y luego reiniciar" es especialmente importante en escenarios de ejecución prolongada, ya que cada token es un recurso valioso.
\end{knowledgebox}

\subsection{Flujo de inicio de sesión}

Al iniciar cada sesión del Coding Agent, se ejecuta un conjunto de procedimientos estandarizados para obtener el estado actual:

\begin{enumerate}
    \item Ejecutar \texttt{pwd} para confirmar el directorio de trabajo
    \item Leer el historial git y los archivos de progreso para comprender el trabajo reciente
    \item Leer la lista de funcionalidades, seleccionando la no completada con mayor prioridad a continuación
    \item Ejecutar \texttt{init.sh} para iniciar el servidor de desarrollo
    \item Realizar pruebas básicas de extremo a extremo para confirmar que las funcionalidades existentes funcionan correctamente
    \item Una vez confirmado que no hay errores de regresión, proceder a implementar la nueva funcionalidad
\end{enumerate}

\begin{warningbox}{Pruebas antes del desarrollo, no desarrollo antes de las pruebas}
Una decisión de diseño clave es: el agente debe verificar primero si las funcionalidades existentes funcionan correctamente antes de comenzar el desarrollo de una nueva. Si se salta este paso y se inicia directamente el desarrollo de la nueva funcionalidad, podría introducir más problemas sobre la base de errores existentes, empeorando la situación. Esto coincide con las mejores prácticas de los ingenieros humanos: después de hacer pull del último código, ejecutar primero un conjunto de pruebas para confirmar que CI está verde antes de comenzar a escribir nuevo código.
\end{warningbox}

\subsection{Secuencia típica de inicio de sesión del agente}

A continuación se muestra la secuencia de comportamiento típico de inicio del Coding Agent:

\begin{lstlisting}[language={}, caption={Flujo típico de inicio de sesión del Coding Agent}]
[Assistant] I'll start by getting my bearings and understanding
            the current state of the project.
[Tool Use]  <bash - pwd>
[Tool Use]  <read - claude-progress.txt>
[Tool Use]  <read - feature_list.json>

[Assistant] Let me check the git log to see recent work.
[Tool Use]  <bash - git log --oneline -20>

[Assistant] Now let me check if there's an init.sh script
            to restart the servers.
<Starts the development server>

[Assistant] Let me verify that fundamental features are
            still working.
<Tests basic functionality>

[Assistant] Core features are working correctly. Let me
            review the feature list to determine what to
            implement next.
<Starts work on a new feature>
\end{lstlisting}

Este flujo de inicio estandarizado reduce significativamente el tiempo de "calentamiento" de cada sesión, evitando que el agente desperdicie tokens en tareas repetitivas como "entender cómo ejecutar el proyecto".

\subsection{Resumen de la sección}

El núcleo del desarrollo incremental es que cada sesión trata solo una funcionalidad y deja el entorno en un estado limpio mediante git commit y archivos de progreso al finalizar la sesión. El flujo de inicio estandarizado de las sesiones garantiza que cada agente pueda comenzar rápidamente, mientras que la estrategia de validar antes de desarrollar previene la acumulación de nuevos problemas sobre los existentes.
